\documentclass[14pt,aspectratio=1610]{beamer}
%\documentclass[14pt,aspectratio=1610,handout]{beamer}
\usepackage{csu-theme}
\usepackage{forest}
\usepackage{booktabs}

\title{\Huge{} Merge Sort}
\subtitle{CSCI 315: \textit{Data Structure Analysis}\\Chapter 18}
%\author{Sean T. Hayes, Ph.D.}
\institute{Charleston Southern University}

%% Comment out date to use default (today)
% \date{February 18, 2025}
\subject{Software Programming}
%\keywords{}

\setlength{\parskip}{0.5em}


\usepackage{linked-lists}
\usepackage{arrays}


\tikzset{
  node of list/.append style = {
    minimum height=6mm,
    minimum width=3mm,
    node distance=6mm,
    font=\small,
  },
  link pointer of list/.append style= {
    minimum height=6mm,
    minimum width=3mm
  },
  pointer/.append style = {
    minimum height=4mm,
    minimum width=4mm,
    node distance=1pt
  },
  label of pointer/.append style = {
    font=\ttfamily\small
  },
  node of array/.append style = {
     minimum height=8mm,
     minimum width=10.8mm,
     font=\large
  },
  label of array el/.append style={
    % rotate=-67.5,
    % anchor=west,
    % align=right,
    % yshift=-3mm,
    font=\sffamily,
    % rotate=-12,
  },
}

% Adapted from the B-Tree slides for Merge Sort diagram
\NewDocumentCommand{\BTreeNode}{ m }{%
	\begingroup
	\foreach\element[count=\i from 1] in {#1} {%
		\ifnum\i>1% After the first element
      {\color{white}\rule[-0.4em]{1.25pt}{1.5em}}%
		\fi%
      \makebox[1.45em][c]{\vspace{1.25em}\element\vspace{1.25em}}%
	}%
	\endgroup%
}

\forestset{
  visible on/.style={
  for tree={
    /tikz/visible on={#1},
    edge+={/tikz/visible on={#1}}}}
}

\tikzset{
  %
  % Used for beamer overlays
  invisible/.style={opacity=0,text opacity=0},
  visible on/.style={alt=#1{}{invisible}},
  alt/.code args={<#1>#2#3}{%
    \alt<#1>{\pgfkeysalso{#2}}{\pgfkeysalso{#3}} % \pgfkeysalso doesn't change the path
  },
}

\begin{document}

\begin{frame}
  \titlepage%
\end{frame}

\section{Merge Sort}
\begin{frame}{Merge Sort}
\begin{itemize}\setlength\itemsep{0.5em}
\item
  \emph{Quick Sort}:\\
    \(O(n\log_2n)\) average case; \(O(n^2)\) worst case
\item
  \emph{Merge Sort}: always \(O(n\log_2n)\)

\item
  Unlike Quick Sort, it divides the list into two sub-lists of nearly equal size.
\item
  Uses the divide-and-conquer technique.
  \begin{itemize}
  \item
    Partitions the list into two sub-lists of nearly equal size.
  \item
    Sorts the sub-lists.
  \item
    Combines the sub-lists into one sorted list.
  \end{itemize}

\end{itemize}
\end{frame}


\begin{frame}{Merge Sort}
\begin{center}
\begin{forest}
  for tree={%
      s sep = 10mm,
      draw, % show border
      edge=-stealth, % show arrow to child
      edge=very thick, % edge weight
      very thick, % node border
      fill=CSUcyan, % node fill
      text=black,% text color
      fit=tight,% options: tight|rectangle|band
      inner ysep=0,
      inner xsep=1pt,
      minimum height=1.5em,
      l=0pt,
      l sep = 10pt,
    }
[,phantom
  [\BTreeNode{35, 28, 17, 45, 62, 48, 30, 38}
    [\BTreeNode{35, 28, 17, 45}
      [\BTreeNode{35, 28}
        [\BTreeNode{35}, alias=div1a]
        [\BTreeNode{28}, alias=div1b]
      ]
      [\BTreeNode{17, 45}
        [\BTreeNode{17}, alias=div2a]
        [\BTreeNode{45}, alias=div2b]
      ]
    ]
    [\BTreeNode{62, 48, 30, 38}
      [\BTreeNode{62, 48}
        [\BTreeNode{62}, alias=div3a]
        [\BTreeNode{48}, alias=div3b]
      ]
      [\BTreeNode{30, 38}
        [\BTreeNode{30}, alias=div4a]
        [\BTreeNode{38}, alias=div4b]
      ]
    ]
  ]
  [\BTreeNode{17, 28, 30, 35, 38, 45, 48, 62}, for tree={grow=north, edge=stealth-, s sep=20.7mm, visible on=<2->}, l=13.75em
    [\BTreeNode{30, 38, 48, 62}
      [\BTreeNode{30, 38}, alias=merge4]
      [\BTreeNode{48, 62}, alias=merge3]
    ]
    [\BTreeNode{17, 28, 35, 45}
      [\BTreeNode{17, 45}, alias=merge2]
      [\BTreeNode{28, 35}, alias=merge1]
    ]
  ]{
    \onslide<2->{
      \draw[-stealth, very thick] (div1a) to (merge1);
      \draw[-stealth, very thick] (div1b) to (merge1);
      \draw[-stealth, very thick] (div2a) to (merge2);
      \draw[-stealth, very thick] (div2b) to (merge2);
      \draw[-stealth, very thick] (div3a) to (merge3);
      \draw[-stealth, very thick] (div3b) to (merge3);
      \draw[-stealth, very thick] (div4a) to (merge4);
      \draw[-stealth, very thick] (div4b) to (merge4);
    }
  }
]
\end{forest}
\end{center}
\end{frame}

\begin{frame}{Merge Sort Algorithm}
  \large
\begin{itemize}\setlength\itemsep{0.5em}
\item
  Uses recursion.
\item
  If the list is of  assize greater than 1,
  \begin{enumerate}
  \item
    Divide the list into two sub-lists.
  \item
    Merge sort the first sub-list.
  \item
    Merge sort the second sub-list.
  \item
    Merge the first sub-list and the second sub-list.
  \end{enumerate}
\end{itemize}
\end{frame}

\section{Merging Arrays}

\begin{frame}{Merging Two Sorted Arrays}
  \begin{tikzpicture}
    \ArrayLabel{}
    \ArrayList{18, 25, 36}
    \invisible{\ArrayNode{I1}}

    \ArrayLabel[0][5]{}
    \ArrayList{20, 30, 44, 46}
    \invisible{\ArrayNode{I2}}

    \alt<-2>{
      \ArrayElLabel[el18]{left}
    } {
      \alt<-8>{
        \ArrayElLabel[el25]{left}
      }{
        \alt<-14>{
          \ArrayElLabel[el36]{left}
        }{
          \ArrayElLabel[elI1]{left}
        }
        
      }
    }
    
    \alt<-5>{
      \ArrayElLabel[el20]{right}
    }{
      \alt<-11>{
        \ArrayElLabel[el30]{right}
      }{
        \alt<-17>{
          \ArrayElLabel[el44]{right}
        }{
          \alt<-20>{
            \ArrayElLabel[el46]{right}
          }{
            \ArrayElLabel[elI2]{right}
          }
        }
      }
    }

    % Temporary Array
    \ArrayLabel[-3]{}

    % [0]
    \alt<1>{
      \ArrayNode{?}
    }{
      \ArrayNode{18}
    }
    \only<-3> {
      \ArrayElLabel{temp}
    }
    
    % [1]
    \alt<-4>{
      \ArrayNode{?}
    }{
      \ArrayNode{20}
    }
    \only<4-6> {
      \ArrayElLabel{temp}
    }

    % [2]
    \alt<-7>{
      \ArrayNode{?}
    }{
      \ArrayNode{25}
    }
    \only<7-9> {
      \ArrayElLabel{temp}
    }

    % [3]
    \alt<-10>{
      \ArrayNode{?}
    }{
      \ArrayNode{30}
    }
    \only<10-12>{
      \ArrayElLabel{temp}
    }

    % [4]
    \alt<-13>{
      \ArrayNode{?}
    }{
      \ArrayNode{36}
    }
    \only<13-15>{
      \ArrayElLabel{temp}
    }

    % [5]
    \alt<-16>{
      \ArrayNode{?}
    }{
      \ArrayNode{44}
    }
    \only<16-18>{
      \ArrayElLabel{temp}
    }

    % [6]
    \alt<-19>{
      \ArrayNode{?}
    }{
      \ArrayNode{46}
    }
    
    \only<19-21>{
      \ArrayElLabel{temp}
    }

    \invisible{\ArrayNode{}}
    \only<22->{
      \ArrayElLabel{temp}
    }
    
    \onslide<22>{}
  \end{tikzpicture}
\end{frame}

\section{Merge Sort on Linked Lists}
\begin{frame}{Merge Sort: Find Middle and Divide}
\begin{tikzpicture}[label of pointer/.append style={rotate=20}, node of list/.append style={inner sep=2pt, minimum width=6.8mm}]
  % \LinkedList{65};
  \LinkPtr[][above]{pHead}
  \LinkedNode{65}
  \LinkedNode{7}
  \LinkedNode{85}
  \LinkedNode{92}
  \LinkedNode{12}
  \alt<-17>{\LinkedNode{9}}{\LinkedNode[null][ele]{9}}
  \LinkedNode{24}
  \LinkedNode{75}
  \LinkedNode{16}

  \only<1>{\TailPtr<-10mm>[null][el12]{pCurr}}
  \only<2>{\TailPtr<-10mm>[el65][el12][out=160, in=330]{pCurr}}
  \only<3-4>{\TailPtr<-10mm>[el7][el12][out=155, in=325]{pCurr}}
  \only<5>{\TailPtr<-10mm>[el85][el12][out=150, in=320]{pCurr}}
  \only<6-7>{\TailPtr<-10mm>[el92][el12][bend left=8]{pCurr}}
  \only<8>{\TailPtr<-10mm>[el12][el12][]{pCurr}}
  \only<9-10>{\TailPtr<-10mm>[el9][el12][bend right=8]{pCurr}}
  \only<11>{\TailPtr<-10mm>[el24][el12][bend right=8]{pCurr}}
  \only<12-13>{\TailPtr<-10mm>[el75][el12][bend right=8]{pCurr}}
  \only<14>{\TailPtr<-10mm>[el16][el12][bend right=8]{pCurr}}
  \onslide<15-16>{\TailPtr<-10mm>[null][el12][bend right=8]{pCurr}}

  \only<-3>{\TailPtr<-10mm>[null][el85]{pMid}}
  \only<4-6>{\TailPtr<-10mm>[el65][el85][out=140, in=320]{pMid}}
  \only<7-9>{\TailPtr<-10mm>[el7][el85][bend left=8]{pMid}}
  \only<10-12>{\TailPtr<-10mm>[el85][el85]{pMid}}
  \only<13-15>{\TailPtr<-10mm>[el92][el85][bend right=8]{pMid}}
  \onslide<16-19>{\TailPtr<-10mm>[el12][el85][bend right=8]{pMid}}

  \only<17->{\TailPtr[el9][el9][]{pFirst2}}
  \only<19->{\TailPtr[el65][el65][]{pFirst1}}
\end{tikzpicture}
% \onslide<8->{%
%   \begin{tikzpicture}[label of pointer/.append style={rotate=20}, link pointer of list/.append style={minimum height=6mm,minimum width=3mm}, node of list/.append style={inner sep=3pt}]
%     \LinkPtr[][above]{pHead}
%     \LinkedNode{65}
%     \LinkedNode{7}
%     \LinkedNode{85}
%     \LinkedNode{92}
%     \LinkedNode{12}
%     \TailPtr{pMid}

%     \LinkPtr[-1.25][above]{pFirst2}
%     \LinkedNode{9}
%     \LinkedNode{24}
%     \LinkedNode{75}
%     \LinkedNode{16}
%   \end{tikzpicture}
% }
\end{frame}

\begin{frame}{Merge Sort: Merge}
Sorted sub-lists are merged into one sorted list.
\begin{itemize}
\item
  Compare elements of sub-lists
\item
  Adjust pointers of nodes with the smaller datum.
\end{itemize}

\begin{tikzpicture}[
    node of list/.append style = {node distance=12.5mm, minimum width=7mm, minimum height = 7.2mm}, 
    %node of list/.append style = {inner sep=4pt},
    %label of pointer/.append style={rotate=22},
    link pointer of list/.append style= { minimum width=5mm, minimum height = 7.2mm }
  ]

  \onslide<2->{\LinkPtr{pHead}}

  \alt<1>{
    \LinkedNode[ptrFlip1][false]{7}
  }{
    \LinkedNode[ptrFlip1]{7}
  }

  \onslide<3-5>{\TailPtr<16mm><-2mm>[ele][last][bend left=10]{pLastMerge}}
  
  \alt<-4>{\LinkedNode{12}}{\LinkedNode[ptrFlip3][false][elFlip2]{12}}
 
  \onslide<9-11>{\TailPtr<15.5mm><-10mm>[last][last][bend left=10]{pLastMerge}}
  \alt<-10>{\LinkedNode{65}}{\LinkedNode[ptrFlip5][false][elFlip4]{65}}


  
  \only<17-19>{\TailPtr<15.5mm><-10mm>[last][last][bend left=10]{pLastMerge}}
  \alt<-18>{\LinkedNode{85}}{\LinkedNode[ptr][false][elFlip5]{85}}

  \only<-3>{      \TailPtr<20mm><0mm>[ el7][el65][out=210, in=30]{pFirst1}}
  \onslide<4-9>{  \TailPtr<20mm><0mm>[el12][el65][bend left=10]{pFirst1}}
  \onslide<10-17>{\TailPtr<20mm><0mm>[el65][el65][bend left=10]{pFirst1}}
  \onslide<18->{  \TailPtr<20mm><0mm>[last][el65][bend left=10]{pFirst1}}

  \LinkedNode{92}

  
  \coordinate (ptr) at (0,-1.5);
  \LinkedNode[ptr][false]{9}
  \onslide<5->{\draw[link] (ptrFlip1) to[bend left=10] (ele);}
  \only<8->{\draw[link] (ptr) to[bend right=20] (elFlip2);}
  \onslide<6-8>{\TailPtr<-4mm><-2mm>{pLastMerge}}

  \alt<-7>{\LinkedNode{16}}{\LinkedNode[ptr][false]{16}}
  
  \onslide<12-13>{\TailPtr<-4mm><-10mm>{pLastMerge}}

  \only<11->{\draw[link] (ptrFlip3) to[bend left=10] (ele);}
  \LinkedNode{24}
  \onslide<14-16>{\TailPtr<-4mm><-10mm>{pLastMerge}}

  \only<16->{\draw[link] (ptr) to[bend right=10] (elFlip4);}
  \alt<-15>{\LinkedNode{75}}{\LinkedNode[ptr][false]{75}}
  \only<19->{\draw[link] (ptrFlip5) to[bend left=10] (el75);}
  \only<22->{\draw[link] (ptr) to[bend right=10] (el85);}

  \only<-6>{      \TailPtr<-8mm><0mm>[ el9][el24][out=130, in=330]{pFirst2}}
  \onslide<7-12>{ \TailPtr<-8mm><0mm>[el16][el24][bend right=10]{pFirst2}}
  \onslide<13-14>{\TailPtr<-8mm><0mm>[el24][el24][bend right=10]{pFirst2}}
  \only<15->{
    \alt<-20>{
      \TailPtr<-8mm><0mm>[ el75][el24][bend left=10]{pFirst2}
    }{
      \TailPtr<-8mm><0mm>[null][el24]{pFirst2}
    }
  }

  \only<20->{\TailPtr<-4mm><8mm>{pLastMerge}}

  \onslide<22>{}
\end{tikzpicture}
\end{frame}

\section{Merge Sort: Analysis}
\begin{frame}{Analysis: Merge Sort}
\begin{itemize}
\item
  Suppose that \(L\) is a list of \(n\) elements, with \(n > 0\).
\item
  Suppose that \(n\) is a power of \(2\); that is, \(n = 2^m\) for \\some integer \(m > 0\), so that we can divide the list \\into two sub-lists, each of size:\\
  \begin{equation*}
    \frac{n}{2}=\frac{2^m}{2}=2^{m-1}
  \end{equation*}
  \(m\) will be the number of recursion levels.
\end{itemize}
\end{frame}

\begin{frame}{Analysis: Merge Sort}
\centering
\begin{forest}
  % before drawing tree={
  %     tikz+={\coordinate (a) at (current bounding box.west);},
  %   },
  % for tree={tier/.option=level},
  % level label/.style={
  %   before typesetting nodes={
  %     for nodewalk={current,tempcounta/.option=level,group={root,tree breadth-first},ancestors}{if={>OR={level}{tempcounta}}{before drawing tree={label me=##1}}{}},
  %   }
  % },
  % label me/.style={tikz+={\node [anchor=base] at (.parent |- a) {#1};}},
  block/.style={, minimum height=1.3cm},
  for tree={
    anchor=center,
    grow=south,
  },
  label/.style={anchor=center,align=left, no edge,font=\footnotesize}
  [,phantom
    [8,tier=l1
      [4,tier=l2
        [2,tier=l3
          [1,tier=l4][1,tier=l4]
        ]
        [2,tier=l3
          [1,tier=l4][1,tier=l4]
        ]
      ]
      [4,tier=l2
        [2,tier=l3
          [1,tier=l4][1,tier=l4]
        ]
        [2,tier=l3
          [1,tier=l4][1,tier=l4]
        ]
      ]
    ]
    [{Recursion Level: 0\\Calls to Merge Sort: 1 with 8 elements each}, label,tier=l1
      [{Recursion Level: 1\\Calls to Merge Sort: 2 with 4 elements each}, label,tier=l2
        [{Recursion Level: 2\\Calls to Merge Sort: 4 with 2 elements each}, label,tier=l3
          [{Recursion Level: 3\\Calls to Merge Sort: 8 with 1 elements each}, label,tier=l4]
        ]
      ]
    ]
  ]
\end{forest}
\end{frame}

\begin{frame}{Analysis: Merge Sort}
\begin{itemize}[<+->]\setlength\itemsep{0.5em}
\item
  To merge two sorted lists of size \(s\) and \(t\), the \\maximum number of comparisons is \(s + t - 1\).
\item
  Merging two sorted lists into a sorted list is where the actual comparisons and assignments are done.
\item
  Max. \# of comparisons at level \(k\) of recursion:
  \begin{equation*}
    2^k\left(\frac{n}{2^k} - 1\right) = n - 2^k = O(n)
  \end{equation*}
\end{itemize}
\end{frame}

\begin{frame}{Analysis: Merge Sort}
\begin{itemize}[<+->]
\item
  There are at most \(O(n)\) comparisons per recursion level.
\item
  Therefore, there are \(O(nm)\) total comparisons, \\where \(m\) is the number of levels of recursion.
\item
  Thus, \(O(nm) \equiv O(n \log_2n)\).

\item
  Key comparisons in the worst case:\\[-1em]
    \begin{equation*}
      O\left(n\log_2n\right)
    \end{equation*}
\item
  Key comparisons in average case:\\[-1em]
    \begin{equation*}
      n\log_2n - 1.25n = O\left(n\log_2n\right)
    \end{equation*}
\end{itemize}
\end{frame}

\begin{frame}{}
\Large\centering{}
\href{https://www.toptal.com/developers/sorting-algorithms}{Sorting Algorithms Animation}
\end{frame}

\setbeamercovered{transparent}
\section{Summary}
\begin{frame}{Summary}
\begin{itemize}[<+->]
\item
  By \emph{asymptotic}, we mean the rate of growth of function \(f(n)\) as \(n\) grows towards infinity.
\item
  On average, a sequential search searches half the list and makes \(O(n)\) comparisons, which is inefficient.
\item
  The binary-search algorithm is the optimal worst-case algorithm for solving search problems by using the \textit{comparison method}.
\item
  A binary search requires the list to be sorted but takes only \(2\log_2n - 3\) or O\((\log_2n)\) key comparisons.
\end{itemize}
\end{frame}


\begin{frame}{Summary}
\begin{itemize}[<+->]
\item
  Bubble Sort: \(O(n^2)\) key comparisons and item assignments.
\item
  Selection Sort: \(O(n^2)\) key comparisons and \(O(n)\) item assignments.
\item
  Insertion Sort: \(O(n^2)\) key comparisons and item assignments.
\item
  Both the Quick Sort and Merge Sort algorithms partition a list to sort it.
  \begin{itemize}
    \item
      Quick Sort: average number of key comparisons is \(O(n\log_2n)\); worst case number of key comparisons is \(O(n^2)\)
    \item
      Merge Sort: the number of key comparisons is \(O(n\log_2n)\) on average and in the worst case.
  \end{itemize}
\end{itemize}
\end{frame}

\section*{Lab 11}
\begin{frame}{Lab 11: Insertion, Quick, and Merge Sort}
\centering\Large
  Let's take a look at Lab 11.
\end{frame}
\end{document}
